\begin{abstract}
\noindent
La boiterie demeure un problème majeur de bien-être et de performance en élevage laitier. Son
repérage repose sur l'observation humaine et le score locomoteur, indispensables, mais difficiles
à maintenir en continu sur un troupeau. Les capteurs d'activité complètent cette surveillance en
signalant des changements comportementaux persistants qui justifient une vérification sur le terrain.
Ce mémoire conçoit et évalue un système reproductible d'\textbf{alerte comportementale précoce}
à partir de capteurs IceQube fixés à la patte.

La contribution principale repose sur la branche temporelle HYPO. Elle compare le comportement
actuel de chaque vache à celui observé pendant sa propre période de référence et recherche une
baisse graduelle, persistante et concordante de
plusieurs familles de signaux: pas, indice de mouvement (\textit{Motion Index}), transitions
posturales, temps couché et temps debout. Cette logique individualisée transforme des séries de données indirectes en priorités de
vérification traçables. L'évaluation repose sur des événements injectés uniquement après la
période de référence, puis comparés à une exécution propre de la même vache afin de ne créditer
que les alertes attribuables. Sur les dégradations graduelles, HYPO recouvre la plupart
des événements et localise mieux les épisodes qu'Isolation Forest et Local Outlier Factor; le
contrôle bref isolé est rejeté. Un test de stress sans nouveau réglage étend l'évaluation à des
départs abrupts, des familles désynchronisées, une récupération asymétrique, du bruit et un
bloc de données manquantes, à aire numérique d'enveloppe appariée. HYPO y garde son avantage sur
IF et LOF et produit moins de fond que le comparateur
pédométrique, restreint au seul canal des pas, sans le dépasser pour la
localisation des formes attendues. Sa cohérence multifamille se manifeste ailleurs: elle
l'empêche de localiser une perturbation limitée aux pas, que le comparateur pédométrique situe
correctement.

Le mémoire explore aussi une extension bidirectionnelle fondée sur une branche nommée INSTABILITÉ,
qui représente des situations où un inconfort locomoteur pourrait s'accompagner d'agitation, de
transitions plus fréquentes ou de fragmentation posturale. Elle reste séparée de HYPO.
L'instabilité isolée produit une sortie de surveillance, tandis que la notification de
vérification demeure réservée aux baisses persistantes ou aux séquences
instabilité-hypoactivité. Son évaluation quantifie le gain de couverture, la charge de
vérification et la sensibilité aux causes non locomotrices.

Une analyse IceTag-SLS synchronisée confronte enfin les sorties aux scores locomoteurs à partir
de données des capteurs strictement antérieures à l'évaluation. Les vaches aux scores les plus
élevés reçoivent davantage de notifications et l'AUC descriptive globale est élevée, mais
l'analyse exacte restreinte au groupe expérimental \textit{Exercise} demeure non concluante:
l'association reste exploratoire, faute d'un effectif suffisant de vaches à score élevé hors de
ce groupe. Le mémoire livre un détecteur temporel interprétable, une interface de priorisation
et un protocole reproductible articulant performance logicielle, hypothèse comportementale et
vérification clinique.

\vspace{1em}
\textbf{Mots-clés:} boiterie bovine, alerte comportementale, capteur IceQube, détection de
changement, élevage de précision.
\end{abstract}

\begin{otherlanguage}{english}
\renewcommand{\abstractname}{Abstract}
\begin{abstract}
\noindent
Lameness remains a major welfare and performance issue in dairy farming. Its detection relies on
human observation and locomotion scoring, essential but difficult to maintain continuously at
herd scale. Activity sensors complement this monitoring by flagging persistent behavioral
changes that justify on-farm verification. This thesis designs and evaluates a reproducible
\textbf{early behavioral-alert} system based on leg-mounted IceQube sensors.

The main contribution is the temporal HYPO branch. It compares each cow with her own baseline
period and searches for a gradual, persistent and concordant decline across several signal
families: steps, motion index, postural transitions, lying time and standing time. This
individualized logic converts indirect behavioral measurements into traceable verification priorities.
Evaluation uses events injected strictly after the baseline period and compares each injected
run with a clean run of the same cow, so that only attributable alerts are credited. On gradual
degradations, HYPO recovers most events and localizes episodes better than Isolation Forest and
Local Outlier Factor; the isolated brief variation, a single short departure from normal behavior, is rejected. A no-retuning stress test
extends the evaluation to abrupt onset, desynchronized families, asymmetric recovery, noise and
a missing-data block with matched numerical perturbation-envelope area. HYPO keeps its advantage
over IF and LOF and yields less background load than the pedometric comparator, without
surpassing it in localizing the expected shapes. Its multi-family coherence shows elsewhere: it
prevents HYPO from localizing a perturbation confined to steps, which the pedometric comparator
locates correctly. Requiring joint evidence across families is what separates this approach from
characterizing perturbations through step counts alone.

The thesis explores a bidirectional extension based on an INSTABILITY branch, which
represents situations in which locomotor discomfort could involve agitation, more frequent
transitions or fragmented posture. It remains separated from HYPO: isolated instability
produces a surveillance output, whereas verification notifications are reserved for persistent
decreases or instability-to-hypoactivity sequences. Its evaluation quantifies the gain in
coverage, the verification load and the sensitivity to non-locomotor causes.

A synchronized IceTag-SLS analysis finally compares the outputs with locomotion scores using
only sensor data collected before the evaluation. Cows with the highest scores receive more
notifications and the overall descriptive AUC is high, but the exact analysis restricted to the
\textit{Exercise} group remains inconclusive: the association stays exploratory, for want of
enough high-scoring cows outside that group. The thesis delivers an interpretable temporal
detector, a prioritization interface and a reproducible protocol connecting software performance,
behavioral hypotheses and clinical verification.

\vspace{1em}
\textbf{Keywords:} bovine lameness, behavioral alert, IceQube sensor, change detection,
precision livestock farming.
\end{abstract}
\end{otherlanguage}
